\scanpage{65}
\mspage{61}
% archive sheet number: 64

\[
B\unsure{|\mathrm{o}|_E|\lvert\circ\rvert_E}C\,,
\]
compte tenu de la \deleted{propriété fondamentale}\inserted{compatibilité \unsure{de}} de $\alpha$ aux disjonctions.
Reste la question
\[
\operatorname{cosupp}^{\unsure{\mathrm{o}|\circ|0}}_{E'}(\overline\alpha(B))\subset\operatorname{supp}^{\mathrm{o}}_{E'}(\overline\alpha(C))
\]
i.e. ceci :

$(***)$ Soit $x'\in E'$, alors $x'|\mathrm{o}|\overline\alpha(B)$ implique $x'\in\operatorname{supp}^{\mathrm{o}}_{E'}(\overline\alpha(C))$, i.e. $\forall y'\in E'$, avec $y'|\mathrm{o}|\overline\alpha(C)$, on a $x'|\mathrm{o}|y'$.

Pour le dire autrement :
\[
x',y'\in E',\quad x'|\mathrm{o}|\overline\alpha(B),\ y'|\mathrm{o}|\overline\alpha(C)\Longrightarrow x'|\mathrm{o}|y'
\]

Il ne semble pas que ceci puisse se déduire de la seule compatibilité de $\alpha$ aux disjonctions, qui affirme \unsure{qque|qqque} chose sur la disjonction \inserted{seulement} d'éléments \inserted{$x',y'$} de $E'$ qui \deleted{sont} correspondent à des éléments donnés $x,y$ de $E$.
% slanted
\margin{Non, en effet, voir plus bas}
\deleted{Il semble donc, pour pouvoir conclure}
D'ailleurs, on a un contre-exemple immédiat, en prenant dans $E$ \inserted{et} $E'$ la ``disjonction tautologique'' $x|\mathrm{o}|y\Longleftrightarrow x\ne y$, d'où
$\unsure{\mathbb\Sigma|\Sigma}_E=\unsure{\mathcal P|\mathscr P|P}(E)$, $\mathbb\Sigma_{E'}=\mathcal P(E')$, et une \uline{application} $\alpha:E\longrightarrow E'$, d'où $\alpha_*:\mathbb\Sigma_E\longrightarrow\unsure{\Sigma|\mathbb\Sigma}_{E'}$ donnée par \deleted{$\unsure{\alpha}$} $A\unsure{\longrightarrow|\longmapsto}\alpha(A)$.
\deleted{Même pour cette \unsure{appli}}
La compatibilité de $\alpha$ aux disjonctions signifie ici $\alpha$ injection, mais $\alpha_*$ n'est compatible \unsure{à|:} $C$
